\chapter{Introdução}\label{cap:introducao}

Na manutenção preditiva, ou baseada em condição, a intervenção é decidida pelo estado medido do equipamento \cite{mobley2002}, e, em máquinas rotativas, a grandeza mais usada para medir esse estado é a vibração. A maior parte dos defeitos mecânicos, como desbalanceamento, desalinhamento, folga e desgaste de rolamentos, altera as forças dinâmicas no interior da máquina, e essas forças se transmitem à carcaça, onde podem ser medidas sem interromper a operação \cite{randall2011livro}. A ISO 20816-3, norma da Organização Internacional de Normalização (ISO) que substituiu normas da série ISO 10816, avalia a vibração de máquinas industriais pelo valor eficaz da velocidade de vibração na carcaça, de \SIrange{10}{1000}{\hertz}, e pela mudança desse valor \cite{iso208163_2022}.

O transdutor de referência nessa medição é o acelerômetro piezoelétrico, que exige condicionamento de sinal dedicado e tem custo alto \cite{albarbar2008}. Um acelerômetro capacitivo de sistemas microeletromecânicos (MEMS) custa de uma a duas ordens de grandeza menos \cite[p.~3]{komarizadehasl2021}. Os modelos de consumo, projetados para orientação e movimento humano, têm banda de algumas centenas de hertz \cite{albarbar2008,fidali2024}.

O custo restringe o monitoramento contínuo de vibração. Segundo \citeonline[p.~7]{garcia2022}, os sistemas de predição baseados em sensores de vibração, ruído e temperatura ficariam caros demais se instalados em todas as máquinas de uma fábrica, porque exigiriam muitos sensores, cabeamento e programação das medições, e por isso são usados só nas máquinas críticas. Para máquinas pequenas, como ventiladores e micromotores, as soluções comerciais de manutenção preditiva são, em geral, superdimensionadas e de custo proibitivo \cite[p.~1]{kolok2025}. Nos preços de varejo sem imposto reunidos por \citeonline[p.~3]{komarizadehasl2021}, só o acelerômetro piezoelétrico custa de \num{380} a \num{1610} euros. Os componentes eletrônicos do transmissor e do \textit{gateway} deste trabalho, cotados no AliExpress, custam R\$~\num{176,04} com o imposto da compra, sem a bateria (Tabela~\ref{tab:custo}).

Numa máquina sem histórico de falhas, a detecção compara cada medição com a operação normal registrada da própria máquina. Na classificação supervisionada, o modelo é treinado com exemplos rotulados de cada condição, a normal e cada tipo de falha. Numa máquina recém-instalada, os exemplos de falha só passam a existir quando os defeitos ocorrem \cite{chandola2009}. Os dados de operação normal existem desde o primeiro dia, e um modelo treinado só com eles, na abordagem semissupervisionada ou de classe única, marca o que se afasta da operação normal \cite{chandola2009,pimentel2014}. Esse modelo indica a existência de uma mudança, e o tipo e a severidade do defeito em geral só são identificados com aprendizado supervisionado \cite[p.~1640]{worden2007}. Neste trabalho, o conjunto de medições de operação normal usado no treino é a linha de base de cada máquina (seção~\ref{sec:fund-deteccao}).


Entre os trabalhos consultados, os nós acadêmicos de baixo custo para manutenção preditiva combinam um microcontrolador da família ESP32, um acelerômetro MEMS e detecção não supervisionada \cite{antonini2023,chevtchenko2023,kolok2025}. \citeonline{antonini2023} treinaram e executaram o Isolation Forest \cite{liu2008} no próprio ESP32, e \citeonline{kolok2025} usaram o mesmo detector num protótipo com ESP32-C6 e o módulo MPU-6050, de consumo. O falso alarme é a fração da operação normal marcada como anomalia. Nos dois nós com Isolation Forest que o informam ou permitem calculá-lo, ele foi de \SI{22}{\percent} a \SI{38}{\percent} \cite{chevtchenko2023,kolok2025}. Os \SI{22}{\percent} são do nó de \citeonline{kolok2025}, cujo artigo não informa a taxa, e foram calculados neste trabalho a partir da acurácia, da precisão e da revocação publicadas (seção~\ref{sec:fund-nos-publicados}).

Os nós consultados informam a acurácia ou o falso alarme em conjuntos de dados próprios \cite{chevtchenko2023,kolok2025}, ou deixam a acurácia fora do escopo \cite{antonini2023}. Nenhum deles avaliou o falso alarme de um detector treinado numa rodada de medição e aplicado a outra rodada da mesma máquina, numa instalação industrial. O protótipo de \citeonline{kolok2025}, o de componentes mais próximos dos deste trabalho, calcula o valor eficaz da aceleração, sem integração para velocidade e sem critério em \si{\milli\meter\per\second}, e não cita a ISO 10816 nem a ISO 20816. Os estudos de adequação de MEMS ao monitoramento de condição testaram sensores com limite superior de \SI{1500}{\hertz} ou mais \cite{albarbar2008,fidali2024}, acima da banda de um MEMS de consumo.

\plano{Lacuna a confirmar pela busca registrada, com bases, termos e data. Até o registro, o parágrafo acima vale só para os trabalhos citados na seção~\ref{sec:fund-posicao}.}

Este trabalho avalia a detecção de falhas por afastamento da operação normal em máquinas industriais, com um sistema de baixo custo que adquire e transmite o sinal bruto de vibração. No sistema, um transmissor com microcontrolador ESP32-C3 e módulo MPU-6050, colado à carcaça do motor, adquire janelas de \num{512} amostras por eixo a \SI{1}{\kilo\hertz}, ou \SI{0,512}{\second} de sinal. Ele as envia pela modulação de longo alcance e baixa potência LoRa (\textit{Long Range}) a um \textit{gateway} com ESP32-S3, que as repassa a uma interface de programação de aplicações (API), e a API as grava num banco de dados (seção~\ref{cap:sistema}). O sistema foi projetado para a coleta em campo, com nove requisitos definidos antes dela, e verificado por testes automatizados e por ensaios de ponta a ponta. A detecção é feita no computador, sobre as janelas exportadas do banco. Seis descritores de cada janela, entre eles a velocidade eficaz e a componente na frequência de rotação, alimentam o detector principal, o Isolation Forest, com limiar fixado para \SI{5}{\percent} de falso alarme na própria linha de base (seções~\ref{sec:descritores} e~\ref{sec:detectores}). Em 2 de novembro de 2026, o sistema adquiriu \campoJanelasOperacao\ janelas em \campoEnsaios\ ensaios de \campoMaquinas\ máquinas de uma instalação industrial (dobradeira, guilhotina hidráulica, calandra piramidal e compressor de parafuso), todas em operação normal (seção~\ref{sec:ensaios}). O mesmo detector foi aplicado também a uma base pública com desbalanceamento real, medida numa bancada de laboratório (\citeauthoronline{mey2020}, \citeyear{mey2020}, \citeyear{mey2022dados}).

O objetivo geral é avaliar se um sistema de aquisição de baixo custo, que transmite o sinal bruto de vibração, permite detectar, por afastamento da operação normal, mudanças de vibração do tamanho que a ISO 20816-3 considera significativo, com falso alarme limitado, em máquinas de uma instalação industrial. Os objetivos específicos são quatro: projetar, construir e verificar o sistema de aquisição diante dos requisitos da coleta em campo (seção~\ref{cap:sistema}); adquirir janelas de sinal bruto em máquinas em operação normal e medir a integridade dos dados entregues (seção~\ref{sec:res-aquisicao}); avaliar o detector nesses dados e numa base pública com desbalanceamento real, por hipóteses e critérios registrados antes do cálculo (seções~\ref{sec:hipoteses} e~\ref{cap:resultados}); e examinar, com análises posteriores declaradas como tais, quanto os resultados dependem do detector, da amostra e da forma da assinatura de falha (seção~\ref{sec:res-posteriores}).

A hipótese central é que, com uma linha de base construída só com operação normal, essa cadeia de aquisição de baixo custo detecta mudanças de vibração do tamanho que a norma considera significativo, com falso alarme limitado entre rodadas da mesma máquina. A ISO 20816-3 classifica a vibração em quatro zonas, e a zona B é a da vibração aceitável para operação de longo prazo sem restrição \cite{iso208163_2022}. Na série ISO 10816, uma mudança de \SI{25}{\percent} do limite superior da zona B é considerada significativa \cite{schwirzer2002}. Para máquinas de \SIrange{15}{300}{\kilo\watt} em fundação rígida, esses \SI{25}{\percent} são \SI{0,7}{\milli\meter\per\second}, o tamanho de referência deste trabalho (seção~\ref{sec:fund-norma}).

A hipótese central se desdobra em três hipóteses derivadas, com critérios de corroboração e de refutação registrados em 6 de novembro de 2026, depois das análises exploratórias de 2 e 5 de novembro e antes de qualquer cálculo que as avaliasse (Quadro~\ref{tab:hipoteses}). H1 é corroborada se o detector treinado na primeira rodada da calandra e da guilhotina marcar até \SI{10}{\percent} das janelas da segunda rodada, medida de \campoIntervaloCalandra\ a \campoIntervaloGuilhotina\ depois, com limite superior do intervalo de confiança de até \SI{20}{\percent}, e refutada se marcar mais de \SI{20}{\percent} em uma das máquinas. H2 é corroborada se a menor mudança detectável na velocidade eficaz de \SIrange{10}{184}{\hertz}, diferença entre o quantil de \SI{95}{\percent} e a mediana do ensaio, ficar abaixo de \SI{0,7}{\milli\meter\per\second} em todos os ensaios com ao menos \num{25} janelas. H3 é corroborada se, com uma assinatura de desbalanceamento somada digitalmente a janelas reais de operação normal, metade das janelas for marcada com amplitude abaixo de \SI{0,7}{\milli\meter\per\second} em cada máquina, com correlação de Spearman de \num{0,9} ou mais entre a amplitude e a fração marcada (seção~\ref{sec:injecao}). Uma quarta hipótese, H4, acrescentada por emenda ao registro antes do cálculo, avalia o detector na base pública. Ela é corroborada se até \SI{10}{\percent} das janelas sem desbalanceamento, medidas depois da remontagem do suporte do desbalanceamento, forem marcadas e se ao menos metade das janelas for marcada em cada nível de desbalanceamento (seção~\ref{sec:base-publica}).

Nenhuma das quatro hipóteses foi corroborada. Nos trechos analisados, as \campoJanelas\ janelas chegaram completas ao banco de dados, e, nas máquinas em operação, o valor eficaz da aceleração ficou de \sinalRazaoMin\ a \sinalRazaoMax\ vezes o do sensor em repouso (seção~\ref{sec:res-aquisicao}). H1 ficou \hUmVeredito, porque o detector marcou \hUmCalandraIFDireto~\% das janelas da segunda rodada da calandra, com intervalo de confiança de \SI{95}{\percent} de \hUmCalandraIFDiretoInf~\% a \hUmCalandraIFDiretoSup~\%, entre os \SI{10}{\percent} do critério de corroboração e os \SI{20}{\percent} do de refutação. H2 foi \hDoisVeredito, porque a menor mudança detectável passou de \SI{0,7}{\milli\meter\per\second} nas duas rodadas da calandra, com \hDoisMmdCalandraUm\ e \hDoisMmdCalandraDois~mm/s, e no compressor, com \hDoisMmdCompressorUm~mm/s. H3 foi \hTresVeredito, porque a fração marcada não chegou a \SI{50}{\percent} até \SI{2,8}{\milli\meter\per\second} de assinatura em nenhuma máquina. A maior fração foi de \hTresFracaoMaxCalandra~\%, na calandra, em que o detector já marcava \apUmFracaoSemAssinaturaCalandra~\% das janelas de teste sem assinatura. H4 foi \hQuatroVeredito. Na base pública, o detector marcou de \hQuatroFracaoTres~\% a \hQuatroFracaoQuatro~\% das janelas com desbalanceamento real e também \hQuatroFracaoZero~\% das janelas sem desbalanceamento medidas depois da remontagem do suporte, acima dos \SI{20}{\percent} do critério de refutação.

Seis análises definidas depois dos resultados, e por isso sem o peso de uma hipótese registrada antes de qualquer cálculo, mostraram quanto os resultados dependem do detector e da amostra (seção~\ref{sec:res-posteriores}). Com as sementes de 0 a \num{99} do gerador de números aleatórios do detector, H1 seria corroborada com \apTresVereditoUmCorroborada\ das \apTresSementes, e H4 seria refutada com \apTresVereditoQuatroRefutada. A menor mudança detectável dos ensaios avaliados ficou de \apCincoRazaoMin\ a \apCincoRazaoMax\ vezes a do sensor em repouso, e a variação entre janelas medida nas máquinas reúne a vibração da máquina e a resposta da fixação do sensor, que os dados não separam. Com os descritores alterados pela assinatura de desbalanceamento além da faixa do treino, o escore do Isolation Forest ficou abaixo do limiar em três das quatro máquinas, o que explica parte da refutação de H3. O sistema entregou os dados de que a avaliação dependia, e, nos ensaios de ponta a ponta, cumpriu a medição sem enlace, a retenção dos dados com a API fora do ar e a atualização sem cabo do \textit{gateway} (seção~\ref{sec:sis-verificacao}).

O trabalho detecta afastamento da operação normal e não diagnostica o tipo de falha. Nenhuma máquina da coleta tinha falha conhecida, e, segundo o operador, todas estavam saudáveis. A detecção de falha foi estimada com assinaturas de desbalanceamento somadas a janelas reais e com o desbalanceamento real da base pública, medido numa bancada com motor pequeno e sem carga (seção~\ref{sec:limitacoes}). A coleta cobriu uma instalação, um dia e no máximo duas rodadas por máquina, separadas por \campoIntervaloCalandra\ na calandra e por \campoIntervaloGuilhotina\ na guilhotina. O filtro interno do sensor, de banda nominal de \SI{184}{\hertz}, limita a medição às componentes ligadas à rotação, e os defeitos localizados de rolamento, que excitam ressonâncias estruturais de alguns quilohertz, não foram avaliados (seção~\ref{sec:fund-rolamento}). A velocidade eficaz de \SIrange{10}{184}{\hertz} é um limite inferior da velocidade de banda larga da norma e não classifica as máquinas nas zonas da norma (seção~\ref{sec:fund-norma}).
